\documentclass[12pt,a4paper]{article}
\usepackage[utf8]{inputenc}
\usepackage[T1]{fontenc}
\usepackage[brazil]{babel}
\usepackage{setspace}
\usepackage{geometry}
\usepackage{indentfirst}
\usepackage{graphicx}
\usepackage{titlesec}
\usepackage{fancyhdr}
\usepackage{times}
\usepackage{tocloft}
\usepackage[none]{hyphenat}
\usepackage{xcolor}
\usepackage{booktabs}
\usepackage{array}
\usepackage{tabularx}
\usepackage{colortbl}
\usepackage{enumitem}
\usepackage{amsmath}
\usepackage{amssymb}
\usepackage{amsthm}
\usepackage{listings}
\usepackage{courier}

\geometry{a4paper, left=3cm, right=2cm, top=3cm, bottom=2cm}
\setlength{\headheight}{14.5pt}
\setlength{\parindent}{1.25cm}
\setlength{\parskip}{0.2cm}
\renewcommand{\baselinestretch}{1.5}
\pagestyle{fancy}
\fancyhf{}
\rhead{Universidade de Rio Verde}
\lhead{Desenvolvimento de Software para Concorrência}
\rfoot{\thepage}

\titleformat{\section}{\bfseries\large\raggedright}{\thesection.}{1em}{}
\titleformat{\subsection}{\bfseries}{\thesubsection}{1em}{}
\newcommand{\sectionbreak}{\clearpage}
\renewcommand{\cftsecleader}{\cftdotfill{\cftdotsep}}
\tocloftpagestyle{empty}

% Configuração para listagem de logs e código
\lstset{
    basicstyle=\scriptsize\ttfamily,
    columns=flexible,
    breaklines=true,
    frame=single,
    backgroundcolor=\color{gray!10},
    extendedchars=true,
    showstringspaces=false
}

\begin{document}
\sloppy

% ====================================================================
% CAPA
% ====================================================================
\begin{titlepage}
\begin{center}
\textbf{\large UNIVERSIDADE DE RIO VERDE}\\[1cm]
\textbf{\large Engenharia de Software}\\[1cm]
\textbf{\large Desenvolvimento de Software para Concorrência}\\[1cm]

\vspace{2cm}

\textbf{\LARGE Atividade Prática: Entrega 02}\\[0.5cm]
\textbf{\Large Soluções de Concorrência e Análise Comparativa}\\[3cm]
\end{center}

\begin{flushright}
\textbf{Docente:} Prof. Gustavo Martins Lima \\[0.2cm]
\textbf{Discentes:} Erick Gabriel Alves Barros, \\
Gabriel Corrêa Costa e \\
Gilberto Aparecido Dias Junior \\[0.2cm]
\textbf{Período:} 7° Período
\end{flushright}

\begin{center}
\vfill
Rio Verde (GO), \today
\end{center}
\end{titlepage}

% ====================================================================
\section{Introdução e Contextualização}
% ====================================================================

\noindent Na primeira etapa deste trabalho, implementou-se uma simulação concorrente de transferências bancárias sem proteção de exclusão mútua. O resultado experimental daquela entrega evidenciou a ocorrência de condições de corrida (Lost Updates), ocasionando a corrupção do saldo consolidado do sistema e violando a integridade física dos dados financeiros.

\noindent O objetivo desta segunda etapa é refatorar o sistema aplicando mecanismos avançados de concorrência que corrijam o problema e mantenham a invariabilidade do saldo consolidado, garantindo proteção contra condições de corrida e, concomitantemente, imunidade a deadlocks (impasses).

% ====================================================================
\section{Solução Concorrente Implementada}
% ====================================================================

\noindent A fim de solucionar o problema de concorrência com segurança e eficiência, implementou-se a classe \texttt{SafeTransferService}. O serviço realiza a exclusão mútua nas contas participantes da transação por meio de blocos aninhados \texttt{synchronized} (monitores em Java).

\noindent Para inviabilizar a ocorrência de deadlocks ao realizar transferências simétricas cruzadas (por exemplo, transferências de A para B e de B para A concorrentemente), a aquisição dos locks ocorre de forma estritamente ordenada. A ordem é baseada na ordenação lexicográfica das chaves identificadoras (\texttt{id}) das contas envolvidas:

\begin{lstlisting}[language=Java]
Account firstLock = origin.getId().compareTo(destination.getId()) < 0 
    ? origin : destination;
Account secondLock = firstLock == origin ? destination : origin;

synchronized (firstLock) {
    synchronized (secondLock) {
        executeTransfer(origin, destination, amount);
    }
}
\end{lstlisting}

\noindent Essa abordagem assegura que qualquer thread necessitando dos recursos de A e B irá sempre solicitar o lock de menor identificador primeiro, eliminando o ciclo de espera mútuo.

% ====================================================================
\section{Justificativa do Mecanismo Adotado}
% ====================================================================

\noindent A escolha por blocos \texttt{synchronized} aninhados com ordenação determinística de identificadores é justificada pelos seguintes aspectos técnicos:

\begin{enumerate}
    \item \textbf{Prevenção Total de Deadlock:} Ao impor uma ordenação global na aquisição de locks, quebra-se a condição de espera circular de Coffman. O sistema elimina a necessidade de lógicas complexas de detecção e recuperação de deadlocks em tempo de execução.
    \item \textbf{Alta Granularidade (Fine-Grained Locking):} Diferente de sincronizar o método inteiro de transferência (o que tornaria todas as transferências estritamente sequenciais no sistema todo), o bloqueio ocorre apenas nas referências das duas contas específicas envolvidas na transação. Isso permite transferências paralelas simultâneas entre outras contas independentes.
    \item \textbf{Monitores Internos da JVM:} Blocos \texttt{synchronized} são primitivas nativas e extremamente otimizadas pelo compilador Java e pela JVM. Além disso, os locks são automaticamente liberados em caso de exceções inesperadas dentro do bloco, impedindo o vazamento de travas (lock leaks) que poderiam ocorrer com APIs de trava manual se não estivessem devidamente envolvidas em blocos try-finally.
\end{enumerate}

% ====================================================================
\section{Testes de Stress e Prevenção de Deadlocks}
% ====================================================================

\noindent O sistema foi submetido a rigorosos testes de estresse em ambiente de simulação concorrente:

\begin{itemize}
    \item \textbf{Teste de Stress de Concorrência:} Execução concorrente com 8 threads simultâneas simulando um montante agregado de 1000 transações paralelas aleatórias entre duas contas sob latência de concorrência.
    \item \textbf{Teste de Imunidade a Deadlock:} Duas threads executando transferências cruzadas repetidas infinitamente em sentidos opostos (A para B e B para A) concorrentemente.
\end{itemize}

\noindent A suíte de testes foi estruturada com barreiras de inicialização do tipo \texttt{CountDownLatch} para garantir partida sincronizada de todas as threads, forçando concorrência máxima de hardware.

% ====================================================================
\section{Análise Comparativa dos Resultados}
% ====================================================================

\noindent A tabela comparativa a seguir consolida o comportamento do sistema antes e depois das correções introduzidas:

\begin{table}[h]
\centering
\caption{Comparativo das Métricas de Execução}
\begin{tabular}{lll}
\toprule
\textbf{Métrica} & \textbf{Versão Ingênua (Entrega 1)} & \textbf{Versão Segura (Entrega 2)} \\
\midrule
Threads Simultâneas & 4 & 8 \\
Volume de Transações & 100 & 1000 \\
Saldo Inicial Total & R\$ 2000,00 & R\$ 2000,00 \\
Saldo Final Total & R\$ 2120,00 (Variável) & R\$ 2000,00 (Estável) \\
Divergência de Saldo & -R\$ 120,00 (Divergente) & R\$ 0,00 (Correto) \\
Ocorrência de Deadlock & Não Observado & Imune (Comprovado em Teste) \\
\bottomrule
\end{tabular}
\end{table}

\noindent O log de simulação a seguir demonstra a consistência final obtida na versão segura:

\begin{lstlisting}[caption={Extrato Final de Execução da Simulação Segura}]
[main] === SIMULATION COMPLETED ===
[main] Final Balance Account A: 1000.00
[main] Final Balance Account B: 1000.00
[main] Actual total balance: 2000.00
[main] Invariant Total: 2000.00
[main] Balance divergence (Corruption): 0.00
\end{lstlisting}

\noindent Conforme evidenciado, a divergência de saldo foi reduzida a zero absoluto mesmo sob stress elevado e com o dobro de threads concorrentes, provando que o mecanismo adotado garantiu consistência lógica e física total ao sistema financeiro concorrente.

% ====================================================================
\section{Limitações, Dificuldades e Conclusão}
% ====================================================================

\subsection{Dificuldades Encontradas}
\noindent A principal dificuldade consistiu em modelar a concorrência em ambiente local de forma que as falhas e correções pudessem ser capturadas de forma reprodutível e robusta por meio de testes automatizados, e configurar o redirecionamento adequado de streams de erros (stderr) em ambientes PowerShell, contornando limitações de codificação e arquivos em branco.

\subsection{Conclusão}
\noindent O desenvolvimento deste projeto incremental permitiu fixar empiricamente a física por trás de condições de corrida e lost updates. Ao aplicar ordenação de locks por chaves identificadoras e exclusão mútua fina via monitores, provou-se que é possível alcançar corretude matemática e imunidade a impasses sob concorrência massiva de execução.

\end{document}
